\section{Risks and critical reflection}\label{sec:13}

This section is the honest account, and it is deliberately the least flattering part of the
document. Several of the findings below are not disclosures but \emph{functional holes} (things
that do not work, rather than things that are merely unvalidated), and they are stated first
because they are the ones worth the most to find.

\subsection{Four functional holes, two of which are now closed}\label{sec:13.1}

Four holes, each with its verdict; \textbf{Appendix \ref{sec:I.1}} carries the mechanism behind every one.

\textbf{1. A damper stuck \emph{wide open} is invisible at the deployed setting.} The statistic is
two-sided, and the deployed \emph{guarded} lower arm catches the low-side fault with a safety
consequence: a sensor drifting downward, which under-ventilates an occupied room and betrays
itself by reading below the 420 ppm outdoor level. It cannot catch a damper stuck open, because
``everybody left the room'' explains a room that stopped getting stuffier at least as well and
nothing inside a ten-minute judging window separates them. The unguarded arm sees both and costs more
false confirmations than R-3 allows. \textbf{An observer with no occupancy signal cannot have both,
and the choice was made in R-3's favour.}

\textbf{2. \texttt{serviced()} had no transport, and now it has one. Closed.} The only exit from a
latched case was unreachable in deployment, so a false latch was \textbf{permanent}.
\texttt{ServiceReported} now travels on \texttt{service/\{zone\}} validated against the open case,
and the dashboard exposes
\texttt{POST /api/zones/\allowbreak \{zone\}/cases/\allowbreak \{caseId\}/serviced}. The whole
\texttt{CONFIRMED} → \texttt{COMPENSATING} → \texttt{SERVICED} → \texttt{WARMUP} →
\texttt{HEALTHY} lifecycle is tested end to end. This matters most to §\ref{sec:13.8}, whose argument is that a false positive
costs one technician visit, an argument that was being made against a system where it cost that
visit \emph{and} a permanently wrong work order.

\textbf{3. The ambiguous fault is absent from the model, so the classifier's hardest discrimination
is never tested.} The ambiguity a real building would produce (a partial obstruction in one duct
branch, reading as both a partly stuck damper and a partly blocked filter) \textbf{is not in the
fault model at all}, and the enumeration cannot express it, so a clean confusion matrix would have
been weak evidence rather than strong.

\textbf{4. And worse than untested: the \texttt{BLOCKED\_FILTER} branch is unreachable with this
physics.} This is the finding I least expected and would lead with in the oral; a unit test pins the
outcome as \texttt{UNCLASSIFIED} so it cannot rot. Two structural reasons, neither a bug to patch:
the proportional split only redistributes a fixed total, so no command the observer judges moves all
six residuals the same way; and the quasi-steady generation estimate absorbs the blockage exactly.
True flow low by $(1-\beta)$, believed flow high by the same factor, so only the time constant
differs. The work order still carries the measured plant power excess, which is right for a case the
system cannot prove and is not a working three-way classifier. \textbf{Of the mechanisms this project
set out to name, exactly one is classified: \texttt{STUCK\_DAMPER}}. \texttt{BLOCKED\_FILTER} can
be pointed at but never confirmed, and \texttt{SENSOR\_DRIFT} is never emitted at all, 0 of 76 runs
from a drift ramp and 0 of 76 from a step offset (§\ref{sec:7.3}, §\ref{sec:11.5}).

\subsection{The detector is evaluated against its own generator}\label{sec:13.2}

The simulator produces exactly the faults the monitor looks for, using the same mass balance and the
same flow split. That invalidates the \textbf{detection-performance numbers as transferable
quantities} (R-1's delay, R-3's false-alarm rate and R-2's comfort percentile are conditional on my
own fault definition) and does \emph{not} invalidate the \textbf{architectural claims}, which are
properties of code: that the monitor never commands, that a latched case has one exit, that stale
readings give \texttt{BLIND} with zero fault events, that a confirmed fault changes strategy without
stopping ventilation. A simulator is the right instrument for those, as it is for R-5 and R-6. Closing
the gap needs measured fault data from a real air handler with service records, or at minimum a second
fault model written independently of the detector. Neither was in scope, and the gap is no smaller for
that.

\subsection{Every physical parameter is invented}\label{sec:13.3}

The full list, so it can be judged rather than taken on trust: the 3 m ceiling, from which every
volume is \texttt{area} $\times$ 3.0; the opening law $a_i = p_i A_i$ with authority proportional to
design flow, and the proportional split; $Q_\text{nom}$ = 0.59 m³/s; $P_\text{nom}$ = 400 W from an
assumed 400 Pa at 0.6 efficiency; \texttt{minimumDamperPosition} = 0.15; the \textbf{affinity
exponent 3}, asserted rather than derived; $\beta$, which has no design value at all and is only a
fault-injection magnitude; drift as an uncapped linear ppm/hour ramp; and the whole occupancy
schedule.

Which of them matter, and how, is the useful part. \textbf{First-order on R-1}: ceiling height,
because every time constant is $V/Q$ and $V$ is linear in it, so a 2.7 m ceiling makes all six rooms
10 \% faster; the design-flow rule, because the observer recovers design occupancy from the
authority and drives suspicion against that ceiling; and the schedule, because a room empty most of
the day hides its own fault. \textbf{First-order on R-2}: \texttt{minimumDamperPosition} and a
sizing that leaves $413/411.4 \approx 1.004$ at the normal speed and 43 \% up to $s=1$, which
together decide whether compensation has anything to work with. \textbf{Near-cosmetic}: $P_\text{nom}$, because the classifier tests a \emph{ratio}
against the affinity-law prediction, so 400 W cancels. \textbf{Worse than cosmetic}: the exponent
3, because the generator and the detector share it, so the blocked-filter power test cannot fail
except through $\beta$. Removing this gap needs a commissioning report: a measured fan curve, a
measured pressure drop, a people counter.

One of them was corrected on measurement rather than argued, and it belongs here because the first
choice was wrong. $Q_\text{nom}$ was initially the design total rounded up, which meets the
simultaneous design condition only at $s=1$, the one speed compensation needs to borrow from. Run
at the normal 0.7 that plant delivers 294 l/s against a 411.4 l/s design total, and the consequence
was not subtle: with \textbf{nothing broken} no room could be held below 1000 ppm at design occupancy
and the monitor confirmed a fault in all six, because a uniformly starved building looks exactly like
a faulty one. Sizing from $\Sigma$ design $\div\ 0.7$ separates the two jobs the rating was doing at
once. \textbf{The fix is not free, and this is the honest cost: margin masks faults.}
\textbf{Appendix \ref{sec:I.2}} gives the correction, its measured price (a filter blocking 35 \% no longer
starves the rooms enough for the monitor to confirm it) and the no-fault arm that now holds every
occupied minute at or below 1000 ppm at design occupancy on all five seeds, which is the figure the
earlier sizing scored 0.000 on.
\subsection{The CUSUM margin is thin, R-3 does fail, and the margin was the symptom rather
than the cause}\label{sec:13.4}
The worst healthy seed peaked at 149.1 ppm against $h$ = 150 ppm at $\sigma$ = 30 ppm, over three
hours in one room: a \textbf{0.9 ppm margin}, with the next seed at 91.6 ppm, so it is a tail draw and
not a design property. R-3 asks for the same statistic across six rooms over seven days, far more
samples, and the measurement contradicts the analytic run length: over 20 fault-free simulated weeks
of all six rooms the monitor raises \textbf{0.55 confirmed cases per week against R-3's limit of one},
with \textbf{16 of 20 weeks clean} and a worst week of 6, against 1.30 per week and 9 clean weeks
before ADR-0010, while on the same tuning R-1's stepped p95 is 6 min against a 15 min budget
(§\ref{sec:11.2}). The scored criterion is the worst week, so R-3 fails.

\textbf{But $h$ was never the knob}, and that is the correction this section needs. ADR-0010
attributed the failure to the \emph{reference} the limit is applied to rather than to the limit: both
plausibility guards were being corrected 30 \% of the way onto every reading, so each CUSUM arm was
accumulating a filter error instead of a residual. Correcting the lower arm's half of it halved the
rate with $k$ = 30 ppm and $h$ = 150 ppm \emph{untouched}. So $k$ and $h$ stay where the spike put
them rather than being retuned until R-3 passes: retuning a limit on the experiment that failed it
is the same methodological error §\ref{sec:13.10} names. The upper arm's half of the same defect is measured,
meets R-3 outright, and was declined because it costs three other requirements (§\ref{sec:11.2}, ADR-0010).
\textbf{Appendix \ref{sec:I.3}} gives the full argument, including which week and which arm the residual
failures sit on.

\subsection{R-7 fails at design occupancy, and the failure is R-3's arriving somewhere else}\label{sec:13.5}

Until this revision R-7 was reported as met, full stop, on 270 observed rooms. It was met on the
occupancy regime it had been run in. With every room at its design head count the same experiment
fails \textbf{231 of 270}, and the shape of the failure is the result.

The two failing clauses are \textbf{disjoint}: 215 observed rooms were themselves \emph{accused of a
fault they did not have} and 16 more sat outside the 5 \% airflow band at a worst deviation of
0.058919 (215 plus 16 is the 231), and they attribute in opposite directions, which the
supervisor-dead pass shows. Accusations run 85 of 90 rooms with the supervisor \emph{dead} against
130 of 180 with it live, so they belong to the regime and the monitor rather than to the response to
the fault; airflow breaches run 0 of 90 dead against 16 of 180 live, so that half is the
compensation's own cost and nothing else's.

\textbf{And the accusation is not where it stops.} Of those 215 falsely latched rooms,
\textbf{130 were also taken off the inner loop}: the supervisor published \texttt{DEGRADED} for
them, confiscating the damper of a room with nothing wrong with it. Zero rooms were degraded without
first being latched, so every one of the 130 is downstream of a false confirmation rather than a
separate defect. That figure reads 0 in any copy of \texttt{T-34b-room-isolation.csv} generated
before the supervisor began re-asserting a degraded zone's mode on every cycle; the confiscation was
always happening, and what changed is that it is now announced often enough to be seen.

\textbf{So the larger half of R-7's design-occupancy failure is R-3's false-alarm rate, measured by a
different experiment in a regime R-3's own replay never visits.} That is the second time this report
has had to say it (§\ref{sec:11.2} records the same pattern inside R-2), and \textbf{R-3 is not one failing
requirement among four; it is the one whose failure the other three inherit}. ADR-0009 predicted this
collision and noted that no experiment asserted R-7 in that regime. One does now. Two bounds travel
with the finding and \textbf{Appendix \ref{sec:I.4}} states both: one of the two airflow columns cannot falsify
this requirement at all, and the 5 \% bound's \emph{reference plant} is ambiguous in a way ADR-0008
never settles (an ADR-level question, not a code change).

\textbf{The mechanism that turns one false alarm into six, which this section owed the reader and
until now left in Appendix \ref{sec:F}.} A false \texttt{STUCK\_DAMPER} does not merely add one accusation;
it \emph{pins} that room's believed opening in \texttt{AhuFlowModel}, and the proportional split is
a ratio over all six rooms. Pin one denominator term and the other five rooms are credited with air
they are not receiving, so their predicted concentrations fall below their readings and they confirm
in turn. When the false-alarm rate was 1.30 per week, instrumenting the replay by room, hour, arm
and whether any room was already pinned showed the two halves to be \textbf{exactly even: 13 of the
26 confirmations were raised with no room yet pinned, and 13 were raised while another already
was}. And of the 17 upper-arm confirmations, 11 were cascaded. One seed latched four rooms in
53 minutes; another latched all six (ADR-0010, Appendix \ref{sec:F}).

\textbf{This is the finding that makes R-3 structural rather than statistical}, and it is why
ADR-0010's response was a re-anchor rather than a larger $h$: a confirmation that pins a damper now
re-anchors the other five rooms and resets their arms, because a step change in believed airflow
makes every ppm banked before it a statement about a different building. Aiming at the cascaded half
took the rate from 1.30 to 0.55 and the clean weeks from 9 to 16 while moving no other requirement.
The half that remains is the primary half, which no re-anchor can reach: \textbf{a shared-plant
design couples its own false alarms, and that coupling is a consequence of ADR-0008's option B
rather than a defect in the detector}. Independent rooms would not cascade, and would also have made
R-7 hold by construction and its experiment measure nothing (§\ref{sec:4.7}).

\subsection{The third diagnostic signal is blind after a confirmation, and reordering the
pipeline does not help}\label{sec:13.6}
§\ref{sec:7.3} records that the stuck-damper leaf fires on two signals rather than three, because a probe can
only be requested once a case is latched and no step has been judged at classification time. The
obvious repair is to reorder: latch as usual, withhold the diagnosis, ask for the probe immediately,
and name the fault when the step lands. \textbf{It was built, measured over 176 paired runs on
identical seeds, and not kept.} It does what it was designed to do (detection moves in 0 of 88
pairs and the fault is named on a genuinely judged step in 28 of the 39 runs that defer, against none
before), and \textbf{the settled classification is identical in all 88 pairs}, because the judged
step reads \texttt{MISSING} for a stuck damper, a blocked filter, nine rooms whose dampers are in
perfect health, and the run with nothing injected. A signal that reads the same for every condition
decides nothing.

The cause is a located defect older than the ordering, and \textbf{correcting that baseline makes the
reorder worse rather than better}, which is the finding worth defending, because by the time a
probe can be commanded the supervisor has degraded the zone, raised the fan and trimmed the other
five, so the room's answer is mostly the compensation's. \textbf{Appendix \ref{sec:I.5}} carries the defect, the
counterfactual and the conclusion, which is the one concrete thing this report would hand to a
successor: \textbf{a three-signal stuck-damper diagnosis needs a step judged \emph{before} a case
opens} (the inner loop's own trims, read against a baseline that survives them) rather than a
probe requested after the response to the fault has already begun.

\subsection{The occupancy dependency is mitigated, not solved}\label{sec:13.7}

The plausibility ceiling of §\ref{sec:7.2} keeps an unmodelled arrival from confirming a fault, and the
sensitivity cost is not small: in A145 (20 seats) or A115 (10 seats), a room holding two people has a
ceiling far above anything a stuck damper produces, so the fault stays invisible until the room fills.
That dependency is now measured over the range a drawn weekday visits: stratifying T-24's
closed-loop runs by the head count at the ignored command gives \textbf{0 of 4 confirmations at 0.625
of design occupancy} and a median falling to 31 min at 1.00 (§\ref{sec:11.2}). And the spike does not cover
it at all: its observer used \emph{true} occupancy, so \textbf{the kill-criterion gate was passed by
a detector the system does not deploy}. What is still missing is exactly this: those figures are a
stratification of runs that happened in one room over 0.625–1.000 of design occupancy, not a sweep
holding the ratio as a controlled input over a wider range. Closing the remainder needs either a real
occupancy signal, which BuildSim has no concept of, or that sweep (which is affordable, and was not
done).

\subsection{Latching is the right trade here, and it is a trade}\label{sec:13.8}

Compensation shrinks the residual that found the fault, so without latching the monitor would clear
every fault it handled successfully, the supervisor would stop compensating, and the pair would
oscillate with a broken damper still broken. The active probe is the right re-test but cannot clear a
case either, because a damper being compensated answers a step exactly as a repaired one does. The
asymmetry is therefore deliberate: \textbf{a sticky false alarm is preferred to a self-clearing real
fault.} For ventilation that is correct: a missed stuck damper means weeks of elevated CO₂, while a
false latch costs one technician visit and a room running with trimmed neighbours. It is the
\emph{wrong} trade wherever the degraded mode is itself the hazard or the expense: latching that shut
ventilation off, or paged an engineer nightly, would make the sticky false positive the more expensive
failure and the clearing policy would have to invert. \textbf{And the cost of a false latch is now
exactly what this argument says it is}, which it was not until recently: the missing
\texttt{serviced()} transport made a false latch permanent, so the trade was being defended against a
system that could not pay the price it quoted. The transport exists (§\ref{sec:13.1}), so a false positive now
costs one technician visit, one HTTP request and a room that ran with trimmed neighbours until it was
made.

\subsection{Scope, and the debts taken knowingly}\label{sec:13.9}

\textbf{Legitimate scope.} Six rooms on one floor behind one air handler is not a building, but the
shared fan makes R-7 non-trivial: keeping five rooms unchanged while compensating the sixth is
something the supervisor must actively do. Simulated time through \texttt{ClockPort} is legitimate
too, and what it buys has to be stated precisely because the obvious phrasing is false of half the
system. \textbf{In the offline harness the replay is reproducible bit for bit}, which is exactly what
R-3's false-alarm count needs. \textbf{In the containerised demo profile it is not}, and that is
measured rather than suspected: four runs on identical \texttt{SEED}, \texttt{TIME\_SCALE} and
\texttt{START\_AT} confirmed A115 at simulated 10:16, 10:17, 10:19 and 10:23. The offline number is
the one every requirement is argued from. What simulated time does not exercise either way is
anything that appears only over real duration: scheduler jitter, clock drift, broker backpressure,
memory growth.

\textbf{Debts.} Skipping a DTO layer is justified by lock-step single-author deployment but makes a
domain refactor a wire-format break; having no message replay is similarly contingent, since a monitor
restart in a real deployment loses the residual history it was accumulating and R-5 then rests
entirely on the journal. Four smaller debts of one shape, the same fact written down twice, have
since been \textbf{corrected}, and \textbf{Appendix \ref{sec:I.6}} keeps them on the record because what they
were is worth more than that they are gone, along with the measured side effect worth more than its
fix: a 0.05 change in one probe amplitude moved \textbf{ten of T-34b's design-occupancy rows from
pass to FAIL}, which says the design-occupancy regime is not merely failing but \emph{unstable}.

\subsection{What I would do differently}\label{sec:13.10}

\noindent\begin{minipage}{\linewidth}
\begin{enumerate}
\tightlist
\item \textbf{Write ADR-0008 before the spike, not after.} The spike tuned $k$ = 30 ppm and
  $h$ = 150 ppm on a 210 m³ room at $Q_\text{max}$ = 0.175 m³/s, parameters none of the six real
  rooms carries, and those values entered ADR-0008's parameter table unchanged.
\item \textbf{Make the occupancy-free observer the spike.} The spike used true occupancy, so the
  kill-criterion gate tested a detector the system does not deploy.
\item \textbf{Run the false-alarm experiment at the requirement's horizon \emph{before} tuning.} A
  149.1 ppm peak against $h$ = 150 should have triggered a seven-day six-room replay immediately;
  tuning a threshold on an experiment shorter than the requirement is this project's clearest
  methodological error.
\item \textbf{Put the ambiguous fault in the generator instead of in the critique.} A per-branch
  obstruction is what would make the confusion matrix worth reporting, and the enumeration cannot
  express it.
\item \textbf{Build the evaluation harness before the dashboard.} The dashboard is pleasant and the
  harness is the evidence.
\item \textbf{Let the supervisor decline to compensate when the compensation is not worth its
  energy.} §\ref{sec:11.7} measures the trade and it is a bad one at design occupancy: a mean of 1.59 times
  the fan power, up to 2.92, to move occupied comfort minutes from 0.0094 to 0.0217 against a
  criterion of 0.95, while on the drawn weekday the plan never raises the fan at all. The
  machinery for refusing already exists: the stuck-shut case reports ``no air path'' and turns the
  work order to occupancy advice. This is the first change I would make next, and it is not made
  here because it would move R-2's measured arm after the result was known.
\end{enumerate}
\end{minipage}

\subsection{Risk register}\label{sec:13.11}

The register is \textbf{Appendix \ref{sec:J}}: each row a risk or an assumption, its impact, and the response
with its status. Three rows now read ``measured, and it does happen'' rather than ``mitigated''
(R-3 at its seven-day six-room horizon, the occupancy-dependent detection delay, and the trim lever
already exhausted before the supervisor is asked for anything), which is the honest shape of this
project's risk position.\label{sec:13end}
